\documentclass[11pt,a4paper]{article}

\usepackage[T1]{fontenc}
\usepackage[a4paper,margin=26mm,headheight=14pt]{geometry}
\usepackage{amsmath,amssymb,amsthm,mathtools,bm,mathrsfs}
\usepackage{newtxtext,newtxmath}
\usepackage{booktabs,array,enumitem}
\usepackage{needspace}
\usepackage{float}
\usepackage{xcolor}
\usepackage{fancyhdr}
\usepackage{microtype}
\usepackage{graphicx}
\usepackage{tikz}
\usetikzlibrary{arrows.meta,positioning,shapes.geometric}
\usepackage{hyperref}
\usepackage[nameinlink,noabbrev]{cleveref}
\newcommand{\needrev}[1]{\textcolor{brightgreen}{#1}}
\crefname{theorem}{theorem}{theorems}
\Crefname{theorem}{Theorem}{Theorems}
\crefname{proposition}{proposition}{propositions}
\Crefname{proposition}{Proposition}{Propositions}
\crefname{lemma}{lemma}{lemmas}
\Crefname{lemma}{Lemma}{Lemmas}
\crefname{corollary}{corollary}{corollaries}
\Crefname{corollary}{Corollary}{Corollaries}
\crefname{definition}{definition}{definitions}
\Crefname{definition}{Definition}{Definitions}
\crefname{remark}{remark}{remarks}
\Crefname{remark}{Remark}{Remarks}
\crefname{equation}{equation}{equations}
\Crefname{equation}{Equation}{Equations}
\crefname{figure}{Figure}{Figures}
\Crefname{figure}{Figure}{Figures}
\usepackage{tcolorbox}
\tcbuselibrary{skins,breakable}
\hypersetup{
  colorlinks=true,
  linkcolor=blue!45!black,
  citecolor=green!35!black,
  urlcolor=blue!55!black,
  pdftitle={Constancy of the Wong Matrix for Three-Dimensional Finite-Dimensional Estimation Algebras of Linear Rank Two and Quadratic Rank Zero},
  pdfsubject={Finite-dimensional estimation algebras and the Wong matrix}
}

\pagestyle{fancy}
\fancyhf{}
\fancyhead[L]{Three-Dimensional Rank-Two Estimation Algebras}
\fancyhead[R]{Constancy of the Wong Matrix}
\fancyfoot[C]{\thepage}

\setlength{\parindent}{1.5em}
\setlength{\parskip}{0.18em}
\setlist[itemize]{leftmargin=2.2em,itemsep=0.2em,topsep=0.3em}
\setlist[enumerate]{leftmargin=2.4em,itemsep=0.2em,topsep=0.3em}
\allowdisplaybreaks[3]
\numberwithin{equation}{section}

\newtheoremstyle{journalplain}
  {0.75em}{0.75em}{\itshape}{}{\bfseries}{.}{0.55em}{}
\newtheoremstyle{journaldef}
  {0.75em}{0.75em}{\normalfont}{}{\bfseries}{.}{0.55em}{}
\theoremstyle{journalplain}
\newtheorem{theorem}{Theorem}[section]
\newtheorem{proposition}[theorem]{Proposition}
\newtheorem{lemma}[theorem]{Lemma}
\newtheorem{corollary}[theorem]{Corollary}
\theoremstyle{journaldef}
\newtheorem{definition}[theorem]{Definition}
\newtheorem{remark}[theorem]{Remark}

\DeclareMathOperator{\ad}{ad}
\DeclareMathOperator{\ord}{ord}
\DeclareMathOperator{\Span}{span}
\DeclareMathOperator{\wt}{wt}
\DeclareMathOperator{\inop}{in}
\newcommand{\R}{\mathbb R}
\newcommand{\C}{\mathbb C}
\newcommand{\N}{\mathbb N}
\newcommand{\U}{\mathcal U}
\newcommand{\dd}{\,\mathrm d}
\newcommand{\pd}{\partial}
\newcommand{\eps}{\varepsilon}
\newcommand{\Lie}[2]{\left[#1,#2\right]}
\newcommand{\zetaSym}{\zeta}

\title{\bfseries
Constancy of the Wong Matrix for Three-Dimensional\\
Finite-Dimensional Estimation Algebras}
\author{Songlin Zhou}
\date{\today}

\begin{document}
\maketitle
\vspace{-2.2em}

\begin{abstract}
For a three-dimensional nonlinear filtering system with smooth drift and
observation functions and orthogonal diffusion, we prove that every
finite-dimensional estimation algebra of linear rank two and quadratic rank
zero has a constant Wong matrix.  The proof remains entirely in the
$C^\infty$ category.  Starting from the published function-element theorem
and affine-structure theorems, we first give a self-contained all-order proof
of constancy of the visible entry $\omega_{12}$: a quotient of normally
ordered principal symbols and triangular growth at every order exclude all
nonzero visible slopes.  We then divide the remaining parameters into three
mutually exclusive sectors.  When a hidden slope is nonzero, a smooth Euler
finite-module lemma shows that the normally ordered coefficients of a
finite-dimensional adjoint module are automatically polynomial in the hidden
Euler variable, and that every weight component is obtained by a genuine Lie
word projection furnished by the Chinese remainder theorem.  This yields
$\partial_3^2$ and then a strictly increasing differential-order Lie ladder.
When the hidden slopes vanish but the visible-slope matrix is nonzero, a
rigidity theorem for the potential followed by a weighted associated-graded
argument excludes hidden profiles of every degree.  The argument covers the
rank-one boundary, arbitrary constant tails, and every zero-parameter branch;
it uses neither analyticity, a maximal-rank hypothesis, nor extrapolation from
finitely many orders.
\end{abstract}

\medskip
\noindent\textbf{Keywords:}
Nonlinear filtering; finite-dimensional estimation algebra; Wong matrix;
non-maximal rank; Euler operator.

\smallskip
\noindent\textbf{2020 Mathematics Subject Classification:}
17B66, 93E11, 35K15.


This paper settles the remaining quadratic-rank-zero branch of the original
classification. The precise result is as follows.

\begin{samepage}
\begin{theorem}[Main theorem]\label{thm:main}
Consider
\begin{align}
\dd x(t)&=f(x(t))\,\dd t+g(x(t))\,\dd v(t),
       &&x(t)\in\R^3,\label{eq:filter-x}\\
\dd y(t)&=h(x(t))\,\dd t+\dd w(t),\label{eq:filter-y}
\end{align}
where $f\in C^\infty(\R^3;\R^3)$,
$h\in C^\infty(\R^3;\R^m)$, and $g\colon\R^3\to O(3)$.  Suppose that the
associated estimation algebra $E$ is finite-dimensional, the state dimension
is three, the linear rank is two, and the quadratic rank is zero.  Then its
Wong matrix $\Omega$ is constant.

More explicitly, in the affine normal form used below,
\begin{align}
\omega_{12}&=b_1x_1+b_2x_2+b_0,\label{eq:affine-w12-general}\\
\omega_{13}&=k_1x_1+k_2x_2+k_3x_3+k_0,\label{eq:affine-w13}\\
\omega_{23}&=h_1x_1+h_2x_2+h_3x_3+h_0,\label{eq:affine-w23}
\end{align}
one necessarily has
\[
b_1=b_2=k_1=k_2=k_3 = h_1 =h_2=h_3=0.
\]
Consequently,
\[
\Omega=
\begin{pmatrix}
0&b_0&k_0\\
-b_0&0&h_0\\
-k_0&-h_0&0
\end{pmatrix}.
\]
\end{theorem}
\end{samepage}


\section{Notations and Prelimiries}

\subsection{Notations}

To avoid confusion with the affine Wong coefficients
$h_0,h_1,h_2,h_3$ below, write the components of the observation map as
$h=(q_1,\ldots,q_m)$.  Set
\begin{equation}
D_i=\pd_i-f_i,\qquad
\eta=\sum_{i=1}^{3}(\pd_i f_i+f_i^2)+\sum_{j=1}^{m}q_j^2,
\qquad
L_0=\frac12\sum_{i=1}^{3}D_i^2-\frac12\eta .
\label{eq:D-eta-L0}
\end{equation}
On expansion,
\[
L_0=\frac12\Delta-\sum_i f_i\pd_i-\sum_i\pd_i f_i
       -\frac12\sum_jq_j^2.
\]
The estimation algebra is
\[
E=\operatorname{Lie}_{\R}\langle L_0,q_1,\ldots,q_m\rangle .
\]
Functions are identified with multiplication operators.  Following the
definition in the literature, let
\[
\mathscr L(E)=E\cap\{\text{homogeneous linear polynomials}\};
\]
the linear rank is $\dim_\R\mathscr L(E)$.  The assertion that the quadratic
rank is zero means here that $E$ contains no nonzero polynomial multiplication
operator of degree exactly two.  These are precisely the two rank assumptions
in the main theorem; they include neither maximal rank, analyticity, nor any
additional elimination condition.  The Wong entries are defined by
\begin{equation}
\omega_{ij}=[D_j,D_i]=\pd_i f_j-\pd_j f_i.
\label{eq:wong-def}
\end{equation}
The definitions of the estimation algebra, linear rank, and Wong matrix agree
with Jiao~\cite[Definitions 2.3--2.5, pp.~21--22]{Jiao2022}; quadratic rank
zero is understood according to the convention in the preceding sentence.

Let $\U_r$ denote the space of smooth-coefficient differential operators of
order at most $r$.  Normal ordering takes the form
\[
P=\sum_\alpha a_\alpha(x)\pd^\alpha,
\]
and
\begin{equation}
(a\pd^\alpha)(b\pd^\beta)
=\sum_{\gamma\le\alpha}\binom{\alpha}{\gamma}
a(\pd^\gamma b)\pd^{\alpha-\gamma+\beta}.
\label{eq:normal-product}
\end{equation}
In particular,
\begin{equation}
[D_i,u]=\pd_i u,\qquad
[L_0,u]=\sum_i(\pd_i u)D_i+\frac12\Delta u,
\label{eq:function-brackets}
\end{equation}
and
\begin{equation}
[L_0,D_i]
=\sum_j\omega_{ij}D_j
 +\frac12\sum_j\pd_j\omega_{ij}
 +\frac12\pd_i\eta .
\label{eq:L0Di}
\end{equation}


\subsection{Prelimiries}


Standard definitions of estimation algebras and the Wong matrix may be found in
Ocone~\cite{Ocone1981}, Shi--Yau~\cite{ShiYau2017}, and Chapter~2 of Jiao's
doctoral dissertation~\cite[pp.~21--26]{Jiao2022}.

Shi and Yau proved that the entries of the Wong matrix of a finite-dimensional
estimation algebra of state dimension three and linear rank two are
polynomials of degree at most one
\cite[Theorems 3.4 and 3.10]{ShiYau2017}.  They also proved the following theorem:
\begin{theorem}\cite[Theorem 1.1, p.~2178]{ShiYau2020}.  
    Let E be the finite dimensional estimation algebra of with
state dimension 3 and rank 2. Then the $\Omega$-matrix has linear structure; i.e., all the
$\Omega$-matrix are degree 1 polynomials.
\end{theorem}
Thus the genuinely unresolved case
is precisely the branch of quadratic rank zero treated here. 
Jiao investigated the remaining rank-two case from the viewpoint of the
structure of the drift and the associated potential function.
Jiao further proved that the
potential function admits a restricted quartic form. More precisely,
Theorem~5.4 of \cite{Jiao2022} shows that
\[
\eta(x)=p_4(x)+\phi(x_3),
\]
where
\[
p_4\in\mathbb R[x_1,x_2,x_3],
\qquad
\deg p_4\le4 .
\]
This implies strong algebraic restrictions on the nonlinear drift.

Furthermore, Theorem~5.5 of \cite{Jiao2022} gives an equivalent
description of the drift structure: the drift vector field can be
decomposed as
\[
f=q+\nabla w,
\]
where \(q\) is a quadratic polynomial vector field and \(w\) is a smooth
potential function. Thus, finite-dimensionality together with the rank
assumption forces the drift to have a quadratic-polynomial component
plus a gradient component.

Jiao then imposed two additional assumptions: the quadratic polynomial
drift condition and an elimination condition introduced in Chapter~5.
Under these extra hypotheses, Theorem~5.12 of \cite{Jiao2022} proves that
the corresponding filter is a Yau filter, which means that the Wong
matrix is constant:
\[
\Omega(x)\equiv\Omega_0 .
\]
Therefore, Jiao's result establishes the constancy of the Wong matrix
for a restricted subclass of the rank-two case. The purpose of the
present work is to remove these additional assumptions and prove the
constancy of the Wong matrix in the general quadratic-rank-zero case.


\begin{theorem}[Ocone's function-element theorem]\label{thm:ocone}
If $E$ is a finite-dimensional estimation algebra and a multiplication
operator $\phi$ belongs to $E$, then $\phi$ is a polynomial of degree at most
two.
\end{theorem}

\noindent
This is due to Ocone~\cite{Ocone1981}; see also
Shi--Yau~\cite[Theorem 2.7, p.~4230]{ShiYau2017} and
Jiao~\cite[Theorem 2.4, p.~25]{Jiao2022}.

\begin{theorem}[Polynomiality of highest-order coefficients]\label{thm:top-coeff}
Let $E$ be a finite-dimensional estimation algebra.  If, for some
$\ell\ge0$,
\[
A=\sum_{|\alpha|=\ell+1}a_\alpha(x)D^\alpha
\pmod{\U_\ell}
\]
belongs to $E$, then every highest-order coefficient $a_\alpha$ is a
polynomial.
\end{theorem}

\noindent
This is the standard result restated by
Shi--Yau~\cite[Theorem 2.9, p.~4230]{ShiYau2017}.  We apply it only to genuine
highest-order coefficients and never use it to control the zeroth-order term
of a first-order operator.

\begin{theorem}[Shi--Yau affine structure]\label{thm:shi-yau-affine}
Suppose that $E$ is finite-dimensional, the state dimension is three, and the
linear rank is two.  After two independent linear functions have been
orthonormalized to $x_1,x_2$,
\[
\omega_{12}\in\R[x_1,x_2],\qquad \deg\omega_{12}\le1,
\]
and $\omega_{13},\omega_{23}$ are affine polynomials in
$x_1,x_2,x_3$.
\end{theorem}

\noindent
The first assertion is~\cite[Theorem 3.4, p.~4233]{ShiYau2017}; the latter
two are~\cite[Theorem 3.10, pp.~4243--4245]{ShiYau2017}.

\begin{corollary}[The function-element space]\label{cor:function-space}
Under the assumptions of the main theorem, an initial orthogonal change of
coordinates can be made so that
\begin{equation}
E\cap C^\infty(\R^3)
=\Span_\R\{1,x_1,x_2\}.
\label{eq:function-space}
\end{equation}
Moreover, $x_1,x_2,D_1,D_2,1\in E$.
\end{corollary}

\begin{proof}
Since the linear rank is two, choose an orthonormal basis of
$\mathscr L(E)$ and then make a constant orthogonal change of coordinates so
that $x_1,x_2\in E$.  Hence
\[
D_i=[L_0,x_i],\qquad [D_i,x_i]=1\qquad(i=1,2),
\]
and therefore $D_1,D_2,1\in E$.  On the other hand, by
\cref{thm:ocone}, every function element is a polynomial of degree at most
two, while quadratic rank zero excludes its degree-two part.  It can thus be
written as $c+\ell$, where $c\in\R$ and $\ell$ is a homogeneous linear
polynomial.  Since $1\in E$,
$\ell=(c+\ell)-c\cdot1\in\mathscr L(E)$, and therefore
$\ell\in\Span\{x_1,x_2\}$.  This proves~\eqref{eq:function-space} and the
membership of all the listed elements.


\subsection{Affine parameters, the Bianchi identity, and a gauge representative}

By~\cref{thm:shi-yau-affine}, initially write
\[
\omega_{23}=h_1x_1+h_2x_2+h_3x_3+h_0.
\]
Here $h_0,h_1,h_2,h_3$ are real coefficients of a Wong entry; the observation
components have uniformly been denoted by $q_j$.
Because $\omega$ is an exterior derivative, one always has
\begin{equation}
\pd_3\omega_{12}-\pd_2\omega_{13}+\pd_1\omega_{23}=0.
\label{eq:Bianchi}
\end{equation}
Note that  Since $\omega_{12}$ is independent of $x_3$, $\partial_3 \omega_{12} = 0, \partial_2\omega_{13} = k_2, \partial_3 \omega_{23} = h_1$. It follows that $h_1=k_2$,
which gives~\eqref{eq:affine-w12-general}--\eqref{eq:affine-w23}.

On $\R^3$, two smooth drifts with the same curl differ by a gradient.  Gauge
conjugation
\[
\mathcal G_\Lambda(P)=e^\Lambda Pe^{-\Lambda}
\]
satisfies
\[
\mathcal G_\Lambda(D_i)=D_i-\pd_i\Lambda;
\]
it therefore adds $\nabla\Lambda$ to the drift, fixes every multiplication
operator, and preserves Lie brackets.  We may consequently choose the
representative
\begin{align}
f_1&=0,\notag\\
f_2&=\frac{b_1}{2}x_1^2+b_2x_1x_2+b_0x_1,\label{eq:f2-general}\\
f_3&=\frac{k_1}{2}x_1^2+k_2x_1x_2+k_3x_1x_3+k_0x_1
 +\frac{h_2}{2}x_2^2+h_3x_2x_3+h_0x_2.
\label{eq:f3-general}
\end{align}
This is also the short Poincare argument of Jiao's Theorem~5.5
\cite[pp.~83--85]{Jiao2022}, reproduced here in full.  Gauge conjugation is
used only as an isomorphism of Lie algebras of differential operators; we do
not claim that the conjugated algebra is the estimation algebra of the same
realizable filtering system.

\end{proof}

\begin{comment}
\begin{lemma}[Transfer of structural input to gauge representatives]
\label{lem:gauge-transfer}
For a Lie algebra $E$ of smooth differential operators on $\R^3$, define
\[
\mathscr F(E)
:=
\left\{
u\in C^\infty(\R^3):
\text{$u$, regarded as a multiplication operator, belongs to $E$}
\right\}.
\]

Let $E_{\rm orig}$ be the estimation algebra of the original filtering
system, and let $E'$ be obtained from $E_{\rm orig}$ by a finite composition
of smooth multiplication conjugations and invertible constant affine
changes of coordinates. Then the following statements hold.

\begin{enumerate}[label=\textup{(\roman*)}]
\item Multiplication conjugation fixes every function element pointwise,
whereas an affine coordinate change transports function elements
bijectively by pullback. Consequently, the identity
\[
\mathscr F(E_{\rm orig})
=\Span_\R\{1,x_1,x_2\},
\]
linear rank two, and the absence of polynomial function elements of total
degree exactly two transfer to the current representative.

\item If $P\in E'$ has ordinary differential order
\[
\ord(P)=r\ge1,
\]
then every coefficient of its ordinary principal symbol $\sigma_r(P)$ is a
polynomial. Equivalently, every coefficient of total $D'$-order $r$ in the
current $D'$-normal form is a polynomial. No conclusion about lower-order
coefficients is asserted.

\item Multiplication conjugation leaves every Wong entry unchanged.
Under an invertible constant affine coordinate change, the Wong two-form
is transformed by an invertible constant pullback. Therefore the property
that the Wong matrix is constant is preserved in both directions.
\end{enumerate}
\end{lemma}


\begin{tcolorbox}[title={Proof of (i), Step 1: multiplication conjugation fixes function elements}]
Let
\[
\mathcal G_\Lambda(P)
:=e^\Lambda P e^{-\Lambda},
\qquad
\Lambda\in C^\infty(\R^3).
\]
For every multiplication operator $u$,
\[
\mathcal G_\Lambda(u)
=e^\Lambda u e^{-\Lambda}
=u.
\]
Since $\mathcal G_{-\Lambda}$ is the inverse conjugation, we have
\[
u\in E
\quad\Longleftrightarrow\quad
u\in\mathcal G_\Lambda(E).
\]
Hence
\[
\mathscr F(\mathcal G_\Lambda(E))
=\mathscr F(E).
\]
Thus multiplication conjugation fixes the function-element space
pointwise, rather than merely mapping it to an abstractly isomorphic
vector space.
\end{tcolorbox}

\begin{tcolorbox}[title={Proof of (i), Step 2: function elements under affine pullback}]
Let
\[
\Phi(y)=Ay+c,
\qquad
A\in{\rm GL}_3(\R).
\]
Define
\[
(U_\Phi v)(y):=v(\Phi(y)),
\qquad
\mathcal C_\Phi(P):=U_\Phi P U_\Phi^{-1}.
\]
For a multiplication operator $u$,
\[
\mathcal C_\Phi(u)=u\circ\Phi.
\]
Consequently,
\[
\mathscr F(\mathcal C_\Phi(E))
=
\{u\circ\Phi:u\in\mathscr F(E)\}.
\]

Let $p$ be a nonzero polynomial of total degree $d$, and write
\[
p=p_d+p_{d-1}+\cdots+p_0,
\]
where $p_j$ is homogeneous of degree $j$ and $p_d\neq0$. The highest
homogeneous part of $p\circ\Phi$ is
\[
p_d(Ay).
\]
Since $A$ is invertible,
\[
p_d(Ay)\neq0.
\]
Therefore
\[
\deg(p\circ\Phi)=\deg p.
\]
Applying the same argument to $\Phi^{-1}$ proves that total polynomial
degree is preserved in both directions.

It follows in particular that the property
\[
\text{``there is no polynomial function element of total degree exactly two''}
\]
is invariant under invertible affine coordinate changes.
\end{tcolorbox}

\begin{tcolorbox}[title={Proof of (i), Step 3: preservation of linear rank two}]
Suppose
\[
\mathscr F(E)
=\Span_\R\{1,x_1,x_2\}.
\]
After the affine change $x=Ay+c$,
\[
\mathscr F(\mathcal C_\Phi(E))
=
\Span_\R
\left\{
1,(Ay+c)_1,(Ay+c)_2
\right\}.
\]
The homogeneous linear parts of $(Ay+c)_1$ and $(Ay+c)_2$ are the first
two row forms of $A$. They are linearly independent because $A$ is
invertible.

A translation may add constants to these linear functions. However,
$1$ belongs to the function-element space, so those constants can be
subtracted. After choosing two adapted homogeneous linear coordinates
$y_1,y_2$, we again obtain
\[
\mathscr F(\mathcal C_\Phi(E))
=\Span_\R\{1,y_1,y_2\}.
\]
Thus the linear rank remains two.
\end{tcolorbox}


\begin{tcolorbox}[title={Proof of (ii), Step 1: a smooth gauge preserves the ordinary principal symbol}]
Let
\[
P
=
\sum_{|\alpha|=r}
a_\alpha(x)\pd^\alpha
\pmod{\U_{r-1}},
\qquad r\ge1.
\]
First,
\[
\mathcal G_\Lambda(\pd_i)
=e^\Lambda\pd_i e^{-\Lambda}
=\pd_i-\pd_i\Lambda.
\]
The second term is an operator of order zero. Induction on $|\alpha|$
therefore gives
\[
e^\Lambda\pd^\alpha e^{-\Lambda}
=
\pd^\alpha
\pmod{\U_{|\alpha|-1}}.
\]
In particular, for $|\alpha|=r$,
\[
e^\Lambda\pd^\alpha e^{-\Lambda}
=
\pd^\alpha
\pmod{\U_{r-1}}.
\]
Since every coefficient $a_\alpha$ is a multiplication operator and is
fixed by the gauge,
\[
\mathcal G_\Lambda(P)
=
\sum_{|\alpha|=r}
a_\alpha(x)\pd^\alpha
\pmod{\U_{r-1}}.
\]
Hence
\[
\sigma_r(\mathcal G_\Lambda(P))
=
\sigma_r(P).
\]
Thus smooth multiplication conjugation preserves ordinary differential
order and leaves the ordinary highest-order principal symbol unchanged.
\end{tcolorbox}

\begin{tcolorbox}[title={Proof of (ii), Step 2: affine transformation of the principal symbol}]
Under
\[
x=Ay+c,
\]
the derivatives transform according to
\[
\pd_{x_i}
=
\sum_s(A^{-1})_{si}\pd_{y_s}.
\]
If
\[
p_r(x,\xi):=\sigma_r(P)(x,\xi),
\]
then the transformed principal symbol is
\[
\sigma_r(\mathcal C_\Phi(P))(y,\zeta)
=
p_r(Ay+c,A^{-T}\zeta).
\]
Therefore the transformed principal coefficients are constant linear
combinations of the functions
\[
a_\alpha(Ay+c).
\]
If all $a_\alpha$ are polynomials, then all transformed coefficients are
polynomials. Applying the inverse affine transformation proves the converse.
Thus polynomiality of all ordinary principal coefficients is preserved in
both directions.
\end{tcolorbox}

\begin{tcolorbox}[title={Proof of (ii), Step 3: relation between \(D\)-normal form and the ordinary symbol}]
Let
\[
D_i=\pd_i-f_i.
\]
Since $f_i$ is a multiplication operator, $D_i-\pd_i$ has differential
order zero. For the ordered monomial
\[
D^\alpha
=
D_1^{\alpha_1}D_2^{\alpha_2}D_3^{\alpha_3},
\]
induction on $|\alpha|$ gives
\[
D^\alpha
=
\pd^\alpha
\pmod{\U_{|\alpha|-1}}.
\]
Consequently, if
\[
Q
=
\sum_{|\alpha|=r}
b_\alpha(x)D^\alpha
\pmod{\U_{r-1}},
\]
then
\[
Q
=
\sum_{|\alpha|=r}
b_\alpha(x)\pd^\alpha
\pmod{\U_{r-1}}.
\]
Hence
\[
\sigma_r(Q)(x,\xi)
=
\sum_{|\alpha|=r}
b_\alpha(x)\xi^\alpha.
\]
Thus the highest-order $D$-normal coefficients are exactly the coefficients
of the ordinary principal symbol. This identification does not apply to
lower-order $D$-slots.
\end{tcolorbox}

\begin{tcolorbox}[title={Proof of (ii), Step 4: application of the highest-coefficient theorem}]
Let $T$ denote the finite composition of the gauge conjugations and affine
coordinate changes used to obtain $E'$ from $E_{\rm orig}$. For
$P\in E'$, define
\[
P_{\rm orig}:=T^{-1}P\in E_{\rm orig}.
\]
The preceding steps show that
\[
\ord(P_{\rm orig})=\ord(P)=r.
\]
Write $P_{\rm orig}$ in the original $D$-normal form:
\[
P_{\rm orig}
=
\sum_{|\alpha|=r}
b_\alpha(x)D^\alpha
\pmod{\U_{r-1}}.
\]
Since $E_{\rm orig}$ is a genuine finite-dimensional estimation algebra,
we may apply \cref{thm:top-coeff} to $P_{\rm orig}$. It follows that every
$b_\alpha$ is a polynomial.

By the preceding subproofs, gauge conjugation leaves the principal symbol
unchanged, while an affine coordinate change preserves polynomiality of
all principal coefficients. Therefore all coefficients of
\[
\sigma_r(P)
\]
are polynomials. Equivalently, the coefficients of total order $r$ in the
current $D'$-normal form are polynomials.
\end{tcolorbox}

\begin{tcolorbox}[title={Proof of (iii), Step 1: the Wong entries are gauge invariant}]
Under multiplication conjugation, the transformed drift is
\[
f_i^\Lambda=f_i+\pd_i\Lambda.
\]
Therefore
\begin{align*}
\omega_{ij}^\Lambda
&=
\pd_i f_j^\Lambda-\pd_j f_i^\Lambda\\
&=
\pd_i(f_j+\pd_j\Lambda)
-\pd_j(f_i+\pd_i\Lambda)\\
&=
\omega_{ij}
+\pd_i\pd_j\Lambda-\pd_j\pd_i\Lambda\\
&=
\omega_{ij}.
\end{align*}
Thus multiplication conjugation fixes every Wong entry pointwise.
In particular,
\[
\Omega^\Lambda=\Omega.
\]
Hence the statement that $\Omega$ is constant is exactly gauge invariant.
\end{tcolorbox}

\begin{tcolorbox}[title={Proof of (iii), Step 2: constancy under affine coordinates}]
Regard the Wong matrix as the coefficient matrix of the two-form
\[
\mathcal W
=
\frac12\sum_{i,j}\omega_{ij}(x)\,
\dd x_i\wedge\dd x_j.
\]
For the affine transformation
\[
x=\Phi(y)=Ay+c,
\]
its pullback is
\[
\Phi^\ast\mathcal W
=
\frac12\sum_{r,s}\omega'_{rs}(y)\,
\dd y_r\wedge\dd y_s,
\]
where
\[
\Omega'(y)
=
A^T\Omega(Ay+c)A.
\]
If $\Omega$ is constant, then $\Omega'$ is constant. Conversely,
\[
\Omega(x)
=
A^{-T}
\Omega'\!\left(A^{-1}(x-c)\right)
A^{-1}.
\]
Thus, if $\Omega'$ is constant, then $\Omega$ is constant as well.
Therefore Wong-matrix constancy is invariant under invertible constant
affine coordinate changes.
\end{tcolorbox}


To simplify notation, the current isomorphic representative will henceforth
still be denoted by $E$.  Whenever a published structure theorem is used, it
is understood to be applied first to $E_{\rm orig}$ and then transferred by
the preceding lemma; an arbitrary gauge-conjugated Lie algebra is never
silently treated as a new filtering estimation algebra.

\end{comment}

\section{Proof of constancy of $\omega_{12}$}

\begin{theorem}
    By direct calculation, we have the following elementary results:
\begin{itemize}
\item We have
\begin{align}
 H_0&=[L_0,D_1]-\frac12\pd_2\omega_{12}-\frac12\pd_3\omega_{13}
=\omega_{12}D_2+\omega_{13}D_3+\frac12\eta_1,\label{eq:H0bar}\\
 H_1&=[L_0,D_2]-\frac12\pd_1\omega_{21}-\frac12\pd_3\omega_{23}
=\omega_{21}D_1+\omega_{23}D_3+\frac12\eta_2.\label{eq:H1bar}
\end{align}
\item Set $X_1 = [L_0,H_1]$ and $X_2 = [L_0,H_2]$, we have
\begin{align}
 X_0&=b_1D_1D_2+b_2D_2^2+k_1D_1D_3+k_2D_2D_3+k_3D_3^2\notag\\
&+\left(-\omega_{12}^2-\omega_{13}^2+\frac12\eta_{11}\right)D_1
+\left(-\omega_{13}\omega_{23}+\frac12\eta_{12}\right)D_2\notag\\
&+\left(\omega_{12}\omega_{23}+\frac12\eta_{13}\right)D_3
\pmod{\mathcal U_0},\label{eq:X0-comp}\\
 X_1&=-b_1D_1^2-b_2D_2D_1+k_2D_1D_3+h_2D_2D_3+h_3D_3^2\notag\\
&+\left(-\omega_{23}\omega_{13}+\frac12\eta_{12}\right)D_1
+\left(-\omega_{21}^2-\omega_{23}^2+\frac12\eta_{22}\right)D_2\notag\\
&+\left(-\omega_{21}\omega_{31}+\frac12\eta_{23}\right)D_3
\pmod{\mathcal U_0}.\label{eq:X1-comp}
\end{align}
\item Set $Y_0 = [L_0,X_0]$ and $Y_1 = [L_0,X_1]$, then $Y_0$ and $Y_1$ are second-order operators: 
\begin{align}
Y_0&=
\left(-3k_1\omega_{13}-3b_1\omega_{12}+\frac12\eta_{111}\right)D_1^2\notag\\
&+\left(b_1\omega_{12}-h_2\omega_{13}-2k_2\omega_{23}+\frac12\eta_{122}\right)D_2^2\notag\\
&+\left(-4b_2\omega_{12}-(3k_2+h_1)\omega_{13}-2k_1\omega_{23}+\eta_{112}\right)D_1D_2\notag\\
&+\left(h_3\omega_{12}+k_1\omega_{13}+k_2\omega_{23}+\frac12\eta_{133}\right)D_3^2\notag\\
&+\left((h_1-k_2)\omega_{12}-4k_3\omega_{13}+2b_1\omega_{23}+\eta_{113}\right)D_1D_3\notag\\
&+\left((h_2+k_1)\omega_{12}+(b_1-h_3)\omega_{13}
+3(b_2-k_3)\omega_{23}+\eta_{123}\right)D_2D_3
\pmod{\mathcal U_1}.\label{eq:Y0-comp}
\end{align}
\begin{align}
Y_1&=
\left(
b_2\omega_{12}
-2k_2\omega_{13}
-k_1\omega_{23}
+\frac12\eta_{112}
\right)D_1^2\notag\\
&+\left(
-3b_2\omega_{12}
-3h_2\omega_{23}
+\frac12\eta_{222}
\right)D_2^2\notag\\
&+\left(
-4b_1\omega_{12}
-2h_2\omega_{13}
-4k_2\omega_{23}
+\eta_{122}
\right)D_1D_2\notag\\
&+\left(
-k_3\omega_{12}
+k_2\omega_{13}
+h_2\omega_{23}
+\frac12\eta_{233}
\right)D_3^2\notag\\
&+\left(
-2b_2\omega_{13}
-4h_3\omega_{23}
+\eta_{223}
\right)D_2D_3\notag\\
&+\left(
-(h_2+k_1)\omega_{12}
-3(b_1+h_3)\omega_{13}
-(b_2+k_3)\omega_{23}
+\eta_{123}
\right)D_1D_3 \pmod{\mathcal U_1}.
\end{align}
\end{itemize}
\end{theorem}

\subsection{The visible second-order class and triangular growth}

\begin{definition}\label{def:G}
Let $\mathcal G$ be the set of all elements of $E$ of order at most two.  For
$B\in\mathcal G$, write its second-order part as
\begin{equation}
B=a_{11}D_1^2+a_{22}D_2^2+a_{33}D_3^2+a_{12}D_1D_2
 +a_{13}D_1D_3+a_{23}D_2D_3
\pmod{\mathcal U_1}.
\label{eq:G-form}
\end{equation}
Note that since $B\in E$,
\[
[[B,x_1],x_1]=2a_{11} \in E,
\quad [[B,x_2],x_2]=2a_{22} \in E,
\quad [[B,x_1],x_2]=a_{12} \in E.
\]
It follows from~\eqref{eq:function-space} that
\[
a_{11},a_{22},a_{12}\in\Span_\R\{1,x_1,x_2\}.
\]
\end{definition}


\begin{theorem}\label{Thm3.1}
Let \(B\in\mathcal G\), and write \(B=Q_B+R_B\pmod{\mathcal U_1}\):
\begin{align}
    Q_B &= a_{11}D_1^2 + a_{12}D_1D_2 + a_{22} D_2^2\\
    R_B &= a_{33}D_3^2 + a_{13} D_1D_3 + a_{23} D_2D_3,
\end{align}
Then for \begin{align}
[H_0,B]
&=[\omega_{12}D_2,Q_B]
+[\omega_{12}D_2,R_B]
+[\omega_{13}D_3,Q_B]
+[\omega_{13}D_3,R_B]
\pmod{\mathcal U_1}.
\end{align}
we have the following elementary results:

(1). $[H_0,B] \in \mathcal G$;

(2). All $[\omega_{12}D_2,R_B]$, $[\omega_{13}D_3,Q_B]$, and $[\omega_{13}D_3,R_B]$ contain operator $D_3$;

(3). We have
\begin{align}
    &[\omega_{12}D_2,Q_B]
    = (\omega_{12}\partial_2a_{22} - b_1a_{12}-2b_2a_{22})D_2^2 \\&
    + (\omega_{12}\partial_2a_{11})D_1^2 + (\omega_{12}\partial_2a_{12} - 2b_1a_{11} - b_2a_{12})D_1D_2. \pmod{\mathcal U_1}
\end{align}
Similarily, $[H_1,B] \in \mathcal G$, all $[\omega_{21}D_1,R_B]$, $[\omega_{23}D_3,Q_B]$, and $[\omega_{13}D_3,Q_B]$ contain operator $D_3$, while
\begin{align}
&[\omega_{21}D_1,Q_B]
=
\left(
-\omega_{12}\partial_1a_{11}
+2b_1a_{11}
+b_2a_{12}
\right)D_1^2
\\
&-\omega_{12}\partial_1a_{22}\,D_2^2+
\left(
-\omega_{12}\partial_1a_{12}
+b_1a_{12}
+2b_2a_{22}
\right)D_1D_2
\pmod{\mathcal U_1}.
\end{align}
\end{theorem}

\begin{tcolorbox}[
  title={Proof of Theorem~\ref{Thm3.1}},
  breakable
]
\begin{proof}
We only prove the assertion for $H_0$. Put
\(
w=\omega_{12},
u=\omega_{13}.
\)
Thus
\[
\omega_{21}=-w,
\qquad
\partial_1w=b_1,\qquad
\partial_2w=b_2,\qquad
\partial_3w=0.
\]
Moreover, since \(B\in\mathcal G\), Definition~\ref{def:G} gives
\[
a_{11},a_{12},a_{22}
\in\operatorname{Span}_{\mathbb R}\{1,x_1,x_2\},
\]
and hence
\[
\partial_3a_{11}
=\partial_3a_{12}
=\partial_3a_{22}=0.
\]

\medskip
\noindent\emph{Proof of the first assertion.}
The order filtration satisfies
\[
[\mathcal U_r,\mathcal U_s]
\subseteq\mathcal U_{r+s-1}.
\]
By construction, \(H_0,H_1\in E\cap\mathcal U_1\), while
\(B\in E\cap\mathcal U_2\). Since \(E\) is a Lie algebra,
\[
[H_i,B]\in E\cap\mathcal U_2=\mathcal G,
\qquad i=0,1.
\]

Furthermore, for some \(\phi_0,\phi_1\in\mathcal U_0\),
\[
H_0=wD_2+uD_3+\phi_0,
\qquad
H_1=-wD_1+vD_3+\phi_1.
\]
Since \(B=Q_B+R_B\pmod{\mathcal U_1}\), the filtration identities
\[
[\mathcal U_0,\mathcal U_2]\subseteq\mathcal U_1,
\qquad
[\mathcal U_1,\mathcal U_1]\subseteq\mathcal U_1
\]
give
\[
\begin{aligned}
[H_0,B]
&=
[wD_2,Q_B]+[wD_2,R_B]\\
&+[uD_3,Q_B]+[uD_3,R_B]
\pmod{\mathcal U_1},
\end{aligned}
\]
and similarly
\[
\begin{aligned}
[H_1,B]
&=
[-wD_1,Q_B]+[-wD_1,R_B]\\
&+[vD_3,Q_B]+[vD_3,R_B]
\pmod{\mathcal U_1}.
\end{aligned}
\]

\medskip
For the remaining assertions, we use the elementary identity
\begin{equation}\tag{*}
\begin{aligned}
[fD_\ell,aD_iD_j]
=
f(\partial_\ell a)D_iD_j
-a(\partial_i f)D_\ell D_j
-a(\partial_j f)D_iD_\ell
\pmod{\mathcal U_1}.
\end{aligned}
\end{equation}

\medskip
\noindent\emph{Proof of the second assertion.}
Define
\[
\mathcal I_3^{(2)}
:=
\left\{
c_{13}D_1D_3+c_{23}D_2D_3+c_{33}D_3^2:
c_{ij}\in C^\infty(\mathbb R^3)
\right\}
+\mathcal U_1.
\]
Thus \(P\in\mathcal I_3^{(2)}\) precisely means that every monomial
in the second-order part of \(P\) contains \(D_3\). Applying \((*)\)
to \([wD_2,R_B]\), the only terms in which the \(D_3\)-factor could
disappear are multiplied by \(\partial_3w=0\). Hence
\[
[wD_2,R_B]\in\mathcal I_3^{(2)}.
\]

For \([uD_3,Q_B]\), the transport term in \((*)\) vanishes because
the coefficients of \(Q_B\) are independent of \(x_3\), while each
remaining term contains \(D_3\). Therefore
\[
[uD_3,Q_B]\in\mathcal I_3^{(2)}.
\]
Finally, in \([uD_3,R_B]\), the transport term retains the original
\(D_3\)-factor and each of the other terms introduces a \(D_3\)-factor.
Consequently,
\[
[uD_3,R_B]\in\mathcal I_3^{(2)}.
\]

\medskip
\noindent\emph{Proof of the third assertion.}
Applying \((*)\) term by term gives
\begin{align*}
[wD_2,a_{11}D_1^2]
&=
(w\partial_2a_{11})D_1^2
-2b_1a_{11}D_1D_2,\\
[wD_2,a_{12}D_1D_2]
&=
-b_1a_{12}D_2^2
+(w\partial_2a_{12}-b_2a_{12})D_1D_2,\\
[wD_2,a_{22}D_2^2]
&=
(w\partial_2a_{22}-2b_2a_{22})D_2^2
\pmod{\mathcal U_1}.
\end{align*}
Adding these identities yields
\[
\begin{aligned}
[wD_2,Q_B]
\equiv{}&
(w\partial_2a_{11})D_1^2\\
&+(w\partial_2a_{22}-b_1a_{12}-2b_2a_{22})D_2^2\\
&+(w\partial_2a_{12}-2b_1a_{11}-b_2a_{12})D_1D_2
\pmod{\mathcal U_1}.
\end{aligned}
\]

\end{proof}
\end{tcolorbox}


\subsection{\texorpdfstring{Classification of the $\mathcal G$ coefficients
and final elimination of the slopes of $\omega_{12}$}
{Classification of the G coefficients and final elimination of the slopes of omega12}}


\begin{lemma}
\label{lem:pure-head-Lie}
Let \(C\in\mathcal G\).

\begin{enumerate}[label=\textup{(\roman*)}]
\item If
\(
[D_1^2][C]=Ax_1+C_0,
\)
then \(A=0\).

\item If
\(
[D_2^2][C]=Bx_2+C_0,
\)
then \(B=0\).
\end{enumerate}
\end{lemma}

\begin{tcolorbox}[
  enhanced,
  breakable,
  title={Proof of the pure-head Lie-growth lemma}
]
\begin{proof}
Define
\[
\mathcal T_m^{(1)}
:=
\left\{
\sum_{\substack{|\alpha|=m\\
\alpha_2+\alpha_3\ge1}}
r_\alpha(x)D^\alpha
\right\}.
\tag{2}
\]
Thus
\(
\mathcal U_{m-1}+\mathcal T_m^{(1)}
\)
contains every term of order at most \(m-1\) and every order-\(m\)
monomial containing \(D_2\) or \(D_3\).

Put
\(
u=Ax_1+C_0.
\)
Then
\[
\partial_1u=A,
\qquad
\partial_2u=\partial_3u=0.
\tag{3}
\]

The second-order part of \(C\) has the decomposition
\[
C=uD_1^2+S
\pmod{\mathcal U_1},
\qquad
S\in\mathcal T_2^{(1)}.
\tag{4}
\]

Equations \((1)\)--\((4)\) give, for every \(m\ge1\),
\[
[\mathcal T_m^{(1)},uD_1^2]
\subseteq
\mathcal T_{m+1}^{(1)}+\mathcal U_m,
\tag{5}
\]
\[
[\mathcal T_m^{(1)},\mathcal T_2^{(1)}]
\subseteq
\mathcal T_{m+1}^{(1)}+\mathcal U_m,
\tag{6}
\]
and, for every \(c\in\mathbb R\),
\[
[cD_1^m,\mathcal T_2^{(1)}]
\subseteq
\mathcal T_{m+1}^{(1)}+\mathcal U_m.
\tag{7}
\]

Define the iteration constructure
\[
P_0=[L_0,C],
\qquad
P_{r+1}=[P_r,C],
\tag{9}
\]
equations \((3)\)--\((7)\) give
\[
\begin{aligned}
P_0
&=
\left[\frac12D_1^2,uD_1^2\right]
+R_0
\pmod{\mathcal U_2}\\
&=
A D_1^3+R_0
\pmod{\mathcal U_2},
\qquad
R_0\in\mathcal T_3^{(1)}.
\end{aligned}
\tag{11}
\]

Assume
\[
P_r
=
c_rD_1^{r+3}+R_r
\pmod{\mathcal U_{r+2}},
\qquad
R_r\in\mathcal T_{r+3}^{(1)}.
\tag{12}
\]
Then
\[
\begin{aligned}
P_{r+1}
={}&
[c_rD_1^{r+3},uD_1^2]
+[c_rD_1^{r+3},S]\\
&+[R_r,uD_1^2]
+[R_r,S]
\pmod{\mathcal U_{r+3}}.
\end{aligned}
\tag{13}
\]
Equations \((5)\)--\((8)\) imply
\[
P_{r+1}
=
(r+3)Ac_rD_1^{r+4}
+R_{r+1}
\pmod{\mathcal U_{r+3}},
\tag{14}
\]
where
\[
R_{r+1}\in\mathcal T_{r+4}^{(1)}.
\tag{15}
\]
Consequently,
\[
c_0=A,
\qquad
c_{r+1}=(r+3)Ac_r,
\tag{16}
\]
and hence
\[
c_r
=
\frac{(r+2)!}{2}A^{r+1}.
\tag{17}
\]

If \(A\ne0\), then the the following iterations $P_{r+1}$ have terms $D_1^{r+3}$ with nonzero coefficients, which is a contradiction to the finite-dimensionality of $E$. Therefore
\(
A=0.
\)

Similarily for the second assertion $B = 0$.
\end{proof}
\end{tcolorbox}



\begin{theorem}\label{prop:G-parallel}
If $(b_1,b_2)\ne(0,0)$, then for every $B\in\mathcal G$ there are real
constants $\lambda,\mu,\gamma$ such that
\begin{equation}
a_{11}=\lambda\omega_{12}+c_1,
\quad a_{22}=\mu\omega_{12}+c_2,
\quad a_{12}=\gamma\omega_{12}+c_{12}.
\label{eq:G-parallel}
\end{equation}
\end{theorem}

\begin{tcolorbox}[
  enhanced,
  breakable,
  title={Proof of Theorem~\ref{prop:G-parallel}}
]
\begin{proof}
Put
\[
w:=\omega_{12}=b_1x_1+b_2x_2+b_0,
\qquad
(b_1,b_2)\ne(0,0).
\tag{30}
\]

\emph{Step 1}: Write
\(
a_{11}
=
\alpha_1x_1+\alpha_2x_2+\alpha_0.
\)
For $B \in \mathcal G$, we define
\(
B_{11}
:=
[\mathsf H_0,B]-b_2B \in \mathcal G.
\)
Using \((35)\),
\[
\begin{aligned}
[D_1^2][B_{11}]
&=
w\partial_2a_{11}-b_2a_{11}\\
&=
\alpha_2w
-b_2(\alpha_1x_1+\alpha_2x_2+\alpha_0)\\
&=
(\alpha_2b_1-\alpha_1b_2)x_1
+
(\alpha_2b_0-b_2\alpha_0).
\end{aligned}
\]
Lemma~\ref{lem:pure-head-Lie}\textup{(i)} gives
\[
\alpha_2b_1-\alpha_1b_2=0.
\]
Therefore
\(
(\alpha_1,\alpha_2)
=
\lambda(b_1,b_2)
\)
for some \(\lambda\in\mathbb R\), and one obtains
\[
\begin{aligned}
a_{11}
&=
\lambda b_1x_1+\lambda b_2x_2+\alpha_0\\
&=
\lambda w+(\alpha_0-\lambda b_0).
\end{aligned}
\]


\emph{Step 2:} Write
\(
a_{22}
=
\beta_1x_1+\beta_2x_2+\beta_0.
\)
Define
\(
B_{22}
:=
[\mathsf H_1,B]+b_1B \in \mathcal G.
\)
Using \((36)\),
\[
\begin{aligned}
[D_2^2][B_{22}]
&=
-w\partial_1a_{22}+b_1a_{22}\\
&=
-\beta_1w
+b_1(\beta_1x_1+\beta_2x_2+\beta_0)\\
&=
(b_1\beta_2-b_2\beta_1)x_2
+
(b_1\beta_0-b_0\beta_1).
\end{aligned}
\]
Lemma~\ref{lem:pure-head-Lie}\textup{(ii)} gives
\[
b_1\beta_2-b_2\beta_1=0.
\]
Therefore
\(
(\beta_1,\beta_2)
=
\mu(b_1,b_2)
\)
for some \(\mu\in\mathbb R\), and one obtains
\[
\begin{aligned}
a_{22}
&=
\mu b_1x_1+\mu b_2x_2+\beta_0\\
&=
\mu w+(\beta_0-\mu b_0).
\end{aligned}
\]


\emph{Step 3:} From \((47)\) and \((56)\),
\[
\partial_1a_{11}=\lambda b_1,
\qquad
\partial_2a_{22}=\mu b_2.
\tag{57}
\]

If
\(
b_2\ne0,
\)
put
\[
C_1:=[\mathsf H_1,B]\in\mathcal G.
\tag{59}
\]
Equation \((47)\), applied to \(C_1\), gives
\[
[D_1^2][C_1]=\Lambda w+C
\tag{60}
\]
for some \(\Lambda,C\in\mathbb R\). Equations \((37)\), \((47)\), and
\((57)\) give
\[
\begin{aligned}
\Lambda w+C
&=
-w(\lambda b_1)
+2b_1(\lambda w+c_1)
+b_2a_{12}\\
&=
\lambda b_1w+2b_1c_1+b_2a_{12}.
\end{aligned}
\tag{61}
\]
Hence
\[
b_2a_{12}
=
(\Lambda-\lambda b_1)w
+
(C-2b_1c_1),
\tag{62}
\]
and therefore
\[
a_{12}
=
\frac{\Lambda-\lambda b_1}{b_2}\,w
+
\frac{C-2b_1c_1}{b_2}.
\tag{63}
\]


If
\(
b_1 \neq 0
\)
we do the same procedure and yield the same result.
\end{proof}
\end{tcolorbox}





\begin{lemma}\label{lem:b-both}
If $\omega_{12}$ is nonconstant, then $b_1b_2\ne 0$.
\end{lemma}

\begin{tcolorbox}[
  title={Exclusion of the one-axis branches},
  breakable
]
\begin{proof}
Assume
\(
b_1=0,b_2\ne0,
\)
and take \(B=\mathsf Y_0\). By Theorem~\ref{prop:G-parallel} and lemma
\(
\mu b_2=0,
\)
hence \(\mu=0.\) $B$ can be formulated as:
\begin{align}
B
&=
(\lambda b_2x_2+c_1)D_1^2
+c_2D_2^2 + 
(\gamma b_2x_2+c_{12})D_1D_2
+R_B
\pmod{\mathcal U_1},
\end{align}

Set
\(
P_0=[L_0,B],
P_{r+1}=[P_r,B],
\)
and define
\(
c_r:=[D_1^{r+1}D_2^2][P_r].
\)
For \(r=0\),
\[
\begin{aligned}
c_0
&=
\partial_2a_{12}+\partial_1a_{22}\\
&=
\partial_2(\gamma b_2x_2+c_{12})
+\partial_1c_2\\
&=
\gamma b_2.
\end{aligned}
\]

Suppose
\[
[D_1^{r+1}D_2^2][P_r]=c_r.
\]
Since
\(
[D_1^{r+1}D_2^2,x_2]
=
2D_1^{r+1}D_2,
\)
one has
\[
\begin{aligned}
&[c_rD_1^{r+1}D_2^2,
  \gamma b_2x_2D_1D_2]=
2\gamma b_2c_rD_1^{r+2}D_2^2
\pmod{\mathcal U_{r+3}}.
\end{aligned}
\]
Hence
\[
c_{r+1}=2\gamma b_2c_r.
\]

The \(\lambda\)-term decreases the \(D_2\)-exponent, the constant terms
produce no new highest-order term, and terms involving \(R_B\) retain a
factor \(D_3\). Thus no other term contributes to the tracked slot.
Therefore
\[
c_r
=
(2\gamma b_2)^rc_0
=
2^r(\gamma b_2)^{r+1},
\]

If \(\gamma\ne0\), then
\(
\operatorname{ord}(P_r)=r+3\to\infty,
\)
contradicting finite-dimensionality. Therefore
\(
\gamma=0.
\)
Thus the \(D_2^2\)- and \(D_1D_2\)-coefficients of \(\mathsf Y_0\) are
constant.

From
\[
[D_2^2][\mathsf Y_0]
=
-h_2\omega_{13}-2k_2\omega_{23}
+\frac12\eta_{122},
\]
differentiation by \(x_1\) gives
\[
0=-h_2k_1-2k_2h_1+\frac12\eta_{1122}.
\]
Using \(h_1=k_2\),
\[
\frac12\eta_{1122}=h_2k_1+2k_2^2.
\]

From
\[
[D_1D_2][\mathsf Y_0]
=
-4b_2\omega_{12}
-(3k_2+h_1)\omega_{13}
-2k_1\omega_{23}
+\eta_{112},
\]
differentiation by \(x_2\) gives
\[
0=-4b_2^2-(3k_2+h_1)k_2-2k_1h_2+\eta_{1122}.
\]
Again using \(h_1=k_2\),
\[
\eta_{1122}
=
4b_2^2+4k_2^2+2k_1h_2.
\]

Doubling the first identity and comparing yields
\(
4b_2^2=0,
\)
contrary to \(b_2\ne0\).

The branch \(b_2=0,\ b_1\ne0\) follows similarily.
\end{proof}
\end{tcolorbox}


\begin{theorem}
\label{thm:w12-constant}
Under the main assumptions,
\[
b_1=b_2=0,
\qquad
\omega_{12}=b_0.
\]
\end{theorem}

\begin{tcolorbox}[
  title={Proof of Theorem~\ref{thm:w12-constant}},
  breakable
]
\begin{proof}
Suppose, to the contrary, that \(\omega_{12}\) is nonconstant. By
\cref{lem:b-both},
\(
b_1b_2\ne0.
\)
Let \(B\in\mathcal G\) be arbitrary. By
\cref{prop:G-parallel}, write $\omega_{12} = w$:
\[
a_{11}=\lambda\omega+c_1,
\qquad
a_{22}=\mu\omega+c_2,
\qquad
a_{12}=\gamma\omega+c_{12}.
\]
For
\(
B'=[\mathsf H_0,B]-b_2B,
\)
\(
[D_1^2][B']=-b_2c_1
\)
is a constant, hence $\lambda' = 0$.
\[
\partial_w\lambda=0,
\qquad
\partial_w\mu=-2\mu b_2-\gamma b_1,
\qquad
\partial_w\gamma=-2\lambda b_1-\gamma b_2.
\]
By lemma,
\[
\partial_w\mu = 0.
\]
If $\partial_w \gamma \neq 0$, then like above we can construct a sequence of iterations $P_r$ such that 
\([D_1^{r+2}D_2][P_r]
=
b_1(\gamma')^{r+1}b_2^r,\) which is a contradiction to the finite-dimensionality. Therefore,
\[
2\mu b_2+\gamma b_1=0,
\qquad
2\lambda b_1+\gamma b_2=0.
\]
Since \(b_1b_2\ne0\), define
\(
\kappa:=\frac{\gamma}{2b_1b_2}.
\)
Then
\[
\gamma=2b_1b_2\kappa,
\mu=-b_1^2\kappa,
\lambda=-b_2^2\kappa.
\]

Now take \(B= Y_0\). The three visible coefficients in
\eqref{eq:Y0-comp} give
\begin{align}
-3k_1\omega_{13}-3b_1\omega_{12}
+\frac12\eta_{111}
&=\lambda\omega_{12}+c_1,
\label{eq:rel1}\\
b_1\omega_{12}-h_2\omega_{13}-2k_2\omega_{23}
+\frac12\eta_{122}
&=\mu\omega_{12}+c_2,
\label{eq:rel2}\\
-4b_2\omega_{12}-(3k_2+h_1)\omega_{13}
-2k_1\omega_{23}+\eta_{112}
&=\gamma\omega_{12}+c_{12}.
\label{eq:rel3}
\end{align}
Differentiate \eqref{eq:rel1} by \(\pd_2\),
\eqref{eq:rel2} by \(\pd_1\), and \eqref{eq:rel3}
by \(\pd_1\) and \(\pd_2\), respectively. This gives
\begin{align}
-3k_1k_2-3b_1b_2+\frac12\eta_{1112}
&=-b_2^3\kappa,
\label{eq:elim1}\\
b_1^2-h_2k_1-2k_2h_1+\frac12\eta_{1122}
&=-b_1^3\kappa,
\label{eq:elim2}\\
-4b_1b_2-(3k_2+h_1)k_1-2k_1h_1+\eta_{1112}
&=2b_1^2b_2\kappa,
\label{eq:elim3}\\
-4b_2^2-(3k_2+h_1)k_2-2k_1h_2+\eta_{1122}
&=2b_1b_2^2\kappa.
\label{eq:elim4}
\end{align}
Using \(h_1=k_2\) and comparing the two pairs containing the same
fourth-order derivatives gives
\[
\kappa=\frac{b_1}{b_1^2+b_2^2},
\qquad
\kappa=
-\frac{b_1^2+2b_2^2}
{b_1(b_1^2+b_2^2)}.
\]
Consequently,
\(
b_1^2+b_2^2=0,
\)
contradicting \(b_1b_2\ne0\). Therefore we prove
\(
b_1=b_2=0,
\omega_{12}=b_0.
\)
\end{proof}
\end{tcolorbox}


\section{Proof of $k_3 = h_3 = 0$}
\subsection{Necessary Lemmas}

\begin{lemma}[Global smooth Euler equation]\label{lem:global-smooth-euler}
Let $Q\in\mathbb C[s]\setminus\{0\}$ and let $A=t\partial_t$. Then
\[
 \ker_{C^\infty(\mathbb R;\mathbb C)}Q(A)
 =\operatorname{span}_{\mathbb C}
   \{t^n:n\in\mathbb N_0,\ Q(n)=0\}.
\]
In particular, every globally smooth solution of an Euler equation is a
finite polynomial.
\end{lemma}

\begin{tcolorbox}[title={Proof of Lemma~\ref{lem:global-smooth-euler}},
  breakable
]
\begin{proof}
For $u\in C^\infty(\mathbb R;\mathbb C)$, the Leibniz rule gives
\[
 \bigl(Q(A)u\bigr)^{(n)}(0)=Q(n)u^{(n)}(0).
\]
Suppose that $Q(A)u=0$, and set
\[
 S=\{n\in\mathbb N_0:Q(n)=0\},
 \qquad
 p(t)=\sum_{n\in S}\frac{u^{(n)}(0)}{n!}t^n.
\]
The set $S$ is finite, $Q(A)p=0$, and $v=u-p$ satisfies $Q(A)v=0$ and is
flat at $t=0$.

We now rigorously exclude nonzero flat solutions. On $t>0$, set $t=e^s$
and $V_+(s)=v(e^s)$; on $t<0$, set $t=-e^s$ and
$V_-(s)=v(-e^s)$. On both half-lines, $A=\partial_s$, and hence
\[
 Q(\partial_s)V_\pm=0.
\]
By ODE theory, if the distinct complex roots of $Q$ are $\mu_j$, with multiplicities $m_j$, the solutions of this constant-coefficient equation are
\[
 V_\pm(s)=\sum_j e^{\mu_js}P_j^\pm(s),
 \qquad \deg P_j^\pm<m_j.
\]
The function $v$ is flat at $t = 0$, and
$v(t) = o(|t|^L)$ for any $L>0$,
\[
 \partial_s^kV_\pm(s)=o(e^{Ls})\qquad(s\to-\infty).
\]
If some $P_j^\pm\neq0$, then we can choose $L>\max_{j:P_j^\pm \neq 0}\operatorname{Re}\mu_j$ contradicts the estimate$O(e^{Ls})$. Hence $V_+=V_-=0$, and consequently $v=0$.
\end{proof}
\end{tcolorbox}

\begin{lemma}[Finite-dimensional Euler decomposition]
\label{lem:smooth-euler-module}
Let $\mathfrak g$ be a finite-dimensional real Lie algebra of
finite-order differential operators on $\mathbb R^2_y\times\mathbb R_t$,
with globally smooth coefficients. Suppose that
\[
J_0=\lambda t\partial_t+r_0(y)\in\mathfrak g,
\qquad
\lambda\in\mathbb R\setminus\{0\},
\qquad
r_0\in C^\infty(\mathbb R^2).
\]
Every $P\in\mathfrak g$ has a unique normal form
\[
P=
\sum_{\alpha\in\mathbb N_0^2,\ q\in\mathbb N_0}
a_{\alpha q}(y,t)\,
\partial_{y_1}^{\alpha_1}
\partial_{y_2}^{\alpha_2}
\partial_t^q.
\]
Then 
\begin{itemize}
\item every coefficient satisfies
\[
a_{\alpha q}(y,t)\in C^\infty(\mathbb R^2)[t].
\]
\item Write
\[
a_{\alpha q}(y,t)
=
\sum_{p\ge0}a_{\alpha qp}(y)t^p
\]
and define
\[
\operatorname{wt}
\left(
a_{\alpha qp}(y)t^p
\partial_{y_1}^{\alpha_1}
\partial_{y_2}^{\alpha_2}
\partial_t^q
\right)
:=p-q.
\]
Then, only finitely many weights occur. If $P_w$ denotes the sum of all
terms of $P$ having weight $w$, then there exists
$\pi_w\in\mathbb R[s]$ such that
\[
P_w=\pi_w(\operatorname{ad}_{J_0})P\in\mathfrak g,
\qquad
\operatorname{ad}_{J_0}(P):=[J_0,P].
\]
Thus every weight component is obtained from actual iterated
Lie brackets with $J_0$.
\end{itemize}
\end{lemma}

\begin{tcolorbox}[
  title={Proof of Lemma~\ref{lem:smooth-euler-module}},
  breakable
]
Define
\[
T(P):=[J_0,P],
\qquad
D(P):=[\lambda t\partial_t,P],
\qquad
N(P):=[r_0(y),P].
\]
Thus
\[
T=D+N.
\]
Since $r_0$ is independent of $t$,
\[
[\lambda t\partial_t,r_0]=0.
\]

For a normal-form term
\(
a(y,t)\partial_y^\alpha\partial_t^q,
\partial_y^\alpha
=
\partial_{y_1}^{\alpha_1}\partial_{y_2}^{\alpha_2},
\)
direct calculation gives
\[
D\left(a\partial_y^\alpha\partial_t^q\right)
=
\lambda(t\partial_t-q)a\,
\partial_y^\alpha\partial_t^q
\]
Hence $D$ preserves the $y$-differential order, whereas every
application of $N$ lowers it by at least one.

Choose a finite basis of $\mathfrak g$. Let $K$ be the largest
$y$-differential order appearing in that basis, then
\[
N^{K+1}P=0
\qquad(P\in\mathfrak g).
\]

Because $J_0\in\mathfrak g$ and $\mathfrak g$ is a Lie algebra, $T = \operatorname{ad}_{J_0}$
is a linear endomorphism of a finite-dimensional vector space.
Let $\chi$ be its characteristic polynomial. By the
Cayley--Hamilton theorem,
\[
\chi(T)P=0
\qquad(P\in\mathfrak g).
\]

Since $D$ and $N$ commute by Jacobi identity, there is a polynomial $R$ in two
commuting variables such that
\[
\chi(D+N)=\chi(D)+NR(D,N).
\]
All operators in this identity commute. Hence
\[
\chi(D)^{K+1}P
=
\bigl(\chi(T)-NR(D,N)\bigr)^{K+1}P= 0.
\]
Indeed, every binomial term containing $\chi(T)$ can be rearranged
so that $\chi(T)$ acts directly on $P$ and therefore vanishes.
The remaining term contains $N^{K+1}$ and also vanishes.

Now write
\[
P=\sum_{\alpha,q}
a_{\alpha q}(y,t)\partial_y^\alpha\partial_t^q.
\]
The operator $D$ does not mix different pairs $(\alpha,q)$.
By uniqueness of normal form,
\[
\left[
\chi\!\left(\lambda(t\partial_t-q)\right)
\right]^{K+1}
a_{\alpha q}(y,t)=0.
\]
For fixed $y$, this is an Euler equation in $t$. The global smooth
Euler lemma therefore gives
\[
a_{\alpha q}(y,t)
=
\sum_{\substack{p\in\mathbb N_0\\
\chi(\lambda(p-q))=0}}
\frac{\partial_t^pa_{\alpha q}(y,0)}{p!}\,t^p \in C^\infty (\mathbb R^2)[t].
\]
Thus $P$ has finite sum
\[P = \sum_{w \in \mathbb Z,\chi(\lambda w) = 0}\sum_{p - q = w} P_w,\]

Note that
\[
D(P_w)=\lambda wP_w.
\]
Moreover, $N$ preserves the weight. Consequently,
\[
(T-\lambda w)^{K+1}P_w=N^{K+1}P_w=0.
\]

The polynomials
\[
(s-\lambda w)^{K+1},
\qquad w\in \{\chi(\lambda w) = 0\},
\]
are pairwise coprime. By the Chinese remainder theorem, for every
$w\in W$ there exists $\pi_w\in\mathbb R[s]$ such that
\[
\pi_w\equiv1\pmod{(s-\lambda w)^{K+1}},
\]
and
\[
\pi_w\equiv0\pmod{(s-\lambda v)^{K+1}}
\qquad(v\in \{w:\chi(\lambda w) = 0\},\ v\ne w).
\]
It follows that
\[
\pi_w(T)P=P_w.
\]
Since $T(P)=[J_0,P]$ preserves $\mathfrak g$,
\[
P_w
=
\pi_w(\operatorname{ad}_{J_0})P
\in\mathfrak g.
\]
\end{tcolorbox}




\subsection{Proof of $k_3 = h_3 = 0$.}


\begin{theorem}\label{thm:sector-I-smooth}
Under the main theorem's assumptions, we have $k_3=h_3=0$.
\end{theorem}

\begin{tcolorbox}[
  title={Proof of Theorem~\ref{thm:sector-I-smooth}},
  breakable
]
Suppose $(k_3,h_3)\neq 0$.
Define the Poincar\'e representative \(f^{\mathrm P}\) by
\begin{align}
f_1^{\mathrm P}&=0,\\
f_2^{\mathrm P}&=b_0x_1,\\
f_3^{\mathrm P}
&=\frac{k_1}{2}x_1^2+k_2x_1x_2+k_3x_1x_3+k_0x_1
 +\frac{h_2}{2}x_2^2+h_3x_2x_3+h_0x_2.
\end{align}
Direct differentiation gives that \(f^{\mathrm P}\) and \(f\) have the same Wong
matrix. The Poincar\'e lemma gives
\(\Phi\in C^\infty(\mathbb R^3)\) such that
\[
f^{\mathrm P}-f=\nabla\Phi.
\]

Define
\[
\mathcal G_\Phi(P):=e^\Phi Pe^{-\Phi},
\qquad
E^{\mathrm P}:=\mathcal G_\Phi(E),
\]
we have
\[
D_i^P
:=\mathcal G_\Phi(D_i)
=\partial_i-f_i^{\mathrm P}.
\]
Furthermore,
\[
L_0^{\mathrm P}
:=\mathcal G_\Phi(L_0)
=\frac12\sum_{i=1}^3(D_i^{\mathrm P})^2
 -\frac12\eta.
\]

The map \(\mathcal G_\Phi\) is a Lie-algebra isomorphism. In particular,
\(D_1^{\mathrm P},D_2^{\mathrm P},L_0^{\mathrm P}\in E^{\mathrm P}\),
and the correspondingly transformed \(H_0^{\mathrm P},H_1^{\mathrm P}\)
belong to \(E^{\mathrm P}\). 

From now on we drop the superscript \({\mathrm P}\) and again write
\(E,D_i,L_0,H_i,\eta\).


Set
\[
\alpha=\frac{k_3k_1+h_3k_2}{\mu},
\qquad
\beta=\frac{k_3k_2+h_3h_2}{\mu},
\qquad
\delta=\frac{k_3k_0+h_3h_0}{\mu},
\]
and
\[
t:=\frac{k_3\omega_{13}+h_3\omega_{23}}{\mu}
=x_3+\alpha x_1+\beta x_2+\delta.
\]
Let
\[
\nabla_1:=\partial_1-\alpha \partial_3,
\qquad
\nabla_2:=\partial_2-\beta \partial_3.
\]
Then
\[
\nabla_1t=\nabla_2t=0,
\qquad
[\nabla_1,\nabla_2]=[\nabla_i,\partial_3]=0.
\]
All normal forms and weights below are taken with respect to this commuting
frame.

The formulas for \(H_0,H_1\) give the genuine element
\[
\begin{aligned}
J
&:=
\frac{k_3(H_0-b_0D_2)+h_3(H_1+b_0D_1)}{\mu}\\
&=
tD_3+\frac{k_3\eta_1+h_3\eta_2}{2\mu}+\frac12
=t\partial_3+r(x_1,x_2,t)
\in E
\end{aligned}
\]
for some smooth function \(r\).
Since \(r-r_0\) vanishes on \(t=0\), Hadamard's lemma gives a smooth
function \(\Lambda\) satisfying
\(
\partial_3\Lambda=\frac{r-r_0}{t}.
\)
For
\[
\widehat E:=e^\Lambda Ee^{-\Lambda},
\]
we therefore have
\[
\widehat J
=e^\Lambda Je^{-\Lambda}
=t\partial_3+r_0(x_1,x_2).
\]

We now apply the smooth Euler finite-module lemma. Since
\[
\widehat D_1
=\nabla_1+\alpha \partial_3-\partial_1\Lambda,
\qquad
\widehat D_2
=\nabla_2+\beta \partial_3-b_0x_1-\partial_2\Lambda,
\]
the Euler finite-module lemma shows that \(\partial_1\Lambda\) and \(\partial_2\Lambda\)
are in $C^\infty[x_1,x_2](t)$.
Every positive-weight component of either function is a multiplication
element of \(\widehat E\). But
\[
\widehat E\cap C^\infty(\mathbb R^3)
=\Span_{\mathbb R}\{1,x_1,x_2\}
\]
has only weight zero. Hence \(\partial_1\Lambda\) and
\(\partial_2\Lambda\) are independent of \(t\). Consequently,
\[
\partial_i(\partial_3\Lambda)=\partial_3(\partial_i\Lambda)=0,
\qquad i=1,2,
\]
and therefore
\[
\partial_3\Lambda=\rho(x_3)
\]
for some \(\rho\in C^\infty(\mathbb R)\).

The Poincar\'e formula for \(f_3\), after substituting
\(
x_3=t-\alpha x_1-\beta x_2-\delta,
\)
has the form
\[
f_3=(k_3x_1+h_3x_2)t+F_0(x_1,x_2).
\]
Direct normal ordering of \(\widehat L_0\) now gives
\begin{align}
\widehat L_0
&=
\frac12\Bigl(
(\nabla_1+\alpha \partial_3)^2
+(\nabla_2+\beta \partial_3)^2
+\partial_3^2
\Bigr)\notag\\
&+a_1(x_1,x_2)\nabla_1+a_2(x_1,x_2)\nabla_2\notag\\
&-\Bigl[
(k_3x_1+h_3x_2)t
+\rho(t-\alpha x_1-\beta x_2-\delta)
+b(x_1,x_2)
\Bigr]\partial_3\notag\\
&+u(x_1,x_2,t),
\label{eq:Lhat-short}
\end{align}
for some smooth \(a_1,a_2,b,u\).

A second application of the Euler finite-module lemma shows that every
coefficient in \eqref{eq:Lhat-short} is polynomial in \(t\). In particular,
setting \(x_1=x_2=0\) in the coefficient of \(\partial_3\) gives
\[
\rho(t-\delta)+b(0,0)\in\mathbb R[t].
\]
Thus \(\rho\) is a one-variable polynomial.

The second-order part of \(\widehat L_0\) is
\[
\frac12\nabla_1^2+\frac12\nabla_2^2
+\alpha\nabla_1\partial_3+\beta\nabla_2\partial_3
+\frac{1+\alpha^2+\beta^2}{2}\partial_3^2.
\]
Its weights are \(0,0,-1,-1,-2\), respectively. Every first-order term
in \eqref{eq:Lhat-short} has weight no smaller than \(-1\), and every term of
\(u\) has nonnegative weight. Hence by the Euler projection lemma
\[
\partial_3^2\in\widehat E.
\]

It remains to obtain unbounded differential order.

(i) If $\deg \rho \ge 2$,
\[
\rho(s)=a_ds^d+\cdots,
\qquad d\ge2,\qquad a_d\neq0.
\]
Then
\(T:
=p_{d-1}(\operatorname{ad}_{\widehat J})(\widehat L_0) = a_dt^d\partial_3 + q(x_1,x_2,t)
\in\widehat E.\)

Assume inductively that \(P_m\in\widehat E\) satisfies
\[
P_m
=
\partial_3^m
\pmod{\mathcal U_{m-1}}.
\]
Set
\[
R_m:=[P_m,[T,P_m]] =
m^2a_d\,d(d-1)t^{d-2}\partial_3^{2m-1}\pmod{\mathcal U_{2m-2}},
\]
we obtain
\[
\begin{aligned}
\operatorname{ad}_{\partial_3^2}^{\,d-2}(R_m)
={}&
2^{d-2}(d-2)!
m^2a_d\,d(d-1)\,
\partial_3^{2m+d-3}\pmod{\mathcal U_{2m+d-4}}.
\end{aligned}
\]
The displayed coefficient is nonzero, which contradicts to the finite-dimentionality.

(ii) If $\deg \rho = 1$, the weight-zero component of
\(\widehat L_0\) is
\[
\begin{aligned}
L^{[0]}:=\Pi_0(\widehat L_0)
={}&
\frac12(\nabla_1^2+\nabla_2^2)
+a_1\nabla_1+a_2\nabla_2\\
&-(k_3x_1+h_3x_2+c)t\partial_3+u_0(x_1,x_2).
\end{aligned}
\]
Assume \(\partial_3^m\in\widehat E\). Then
\[
X_m
:=
\frac{[L^{[0]},\partial_3^m]-mc\partial_3^m}{m}
=
(k_3x_1+h_3x_2)\partial_3^m
\in\widehat E.
\]
A direct calculation gives
\[
[L^{[0]},X_m]
=
\bigl(
k_3\nabla_1+h_3\nabla_2+M_m(x_1,x_2)
\bigr)\partial_3^m,
\]
where
\[
\begin{aligned}
M_m
=
k_3a_1+h_3a_2+
m(k_3x_1+h_3x_2+c)(k_3x_1+h_3x_2)
\end{aligned}
\]
is a multiplication operator. Hence \(M_m\partial_3^m\) commutes with \(X_m\), and
\[
\begin{aligned}
[X_m,[L^{[0]},X_m]]
&=
[k_3x_1+h_3x_2,k_3\nabla_1+h_3\nabla_2]\partial_3^{2m}\\
&=
-(k_3^2+h_3^2)\partial_3^{2m}\\
&=
-\mu \partial_3^{2m}.
\end{aligned}
\]
Since \(\mu>0\), we can derive iteratively $\partial_3^{2m} \in \widehat E$, which contradicts to the finite-dimentionality.
\end{tcolorbox}


\section{Proof of $k_1 = k_2 = h_2 = 0$}


\begin{proof}
Suppose, to the contrary, that
\[
M=
\begin{pmatrix}
k_1&k_2\\
k_2&h_2
\end{pmatrix}
\ne0.
\]
Since \(M\) is real symmetric, an orthogonal change of the visible
variables diagonalizes it. After interchanging the two visible axes if
necessary, and then translating \(x_1\), we may assume
\[
k_2=0,\qquad k:=k_1\ne0,\qquad k_0=0.
\]
Write \(h=h_2\). 

After rechoosing the Poincar\'e representative, we have
\[
f_1=0,\qquad f_2=b_0x_1,\qquad
f_3=F(x_1,x_2),
\]
where
\[
F(x_1,x_2)
=\frac{k}{2}x_1^2+\frac{h}{2}x_2^2+h_0x_2.
\]
Thus
\[
D_1=\partial_1,\qquad
D_2=\partial_2-b_0x_1,\qquad
D_3=\partial_3-F.
\]
All coordinate changes and gauge conjugations used below are Lie-algebra
isomorphisms and therefore preserve finite-dimensionality.

\medskip
\noindent\emph{Step 1: the only possible \(x_3\)-dependence of \(\eta\).}

In the present normal form,
\[
H_0=[L_0,D_1]
=b_0D_2+kx_1D_3+\frac12\eta_1,
\]
and
\[
H_1=[L_0,D_2]
=-b_0D_1+(hx_2+h_0)D_3+\frac12\eta_2.
\]
Direct normal ordering gives
\begin{align*}
[D_1,[D_1,H_0]]
 &=\frac12\partial_1^3(\eta-F^2),\\
[D_2,[D_1,H_0]]
 &=\frac12\partial_1^2\partial_2(\eta-F^2),\\
[D_2,[D_2,H_0]]
 &=\frac12\partial_1\partial_2^2(\eta-F^2),\\
[D_2,[D_2,H_1]]
 &=\frac12\partial_2^3(\eta-F^2).
\end{align*}
The four left-hand sides are multiplication operators in \(E\). Since
\[
E\cap C^\infty(\mathbb R^3)
=\operatorname{span}_{\mathbb R}\{1,x_1,x_2\},
\]
they are independent of \(x_3\). Consequently, for arbitrary \(s,t\),
the function
\[
\eta(x_1,x_2,s)-\eta(x_1,x_2,t)
\]
is a polynomial of degree at most two in \(x_1,x_2\). Therefore. $\eta$ can be written as
\[
\eta
=\eta_0(x_1,x_2)
+(x_1,x_2)^TQ(x_3)(x_1,x_2)+\ell(x_3)\cdot (x_1,x_2)+e(x_3),
\]
where \(Q(x_3)\) is symmetric, \(\ell(x_3)\in\mathbb R^2\), and
\(\eta_0\) is smooth.

Now \([D_2,H_0]\) is a function element, so the off-diagonal entry of
\(Q\) is constant. Moreover,
\[
h[D_1,H_0]-k[D_2,H_1]
\]
is a function element because its \(D_3\)-coefficient is zero. It follows
that
\[
hQ_{11}-kQ_{22}
\]
is constant. Absorbing the constant quadratic terms into \(\eta_0\), we
may therefore write
\[
Q(x_3)
=\varepsilon(x_3)
\begin{pmatrix}
k&0\\
0&h
\end{pmatrix}
\]
for some smooth function \(\varepsilon\).

The genuine Lie element
\[
T:=\frac1k[D_1,H_0]
\]
then has the form
\[
T=D_3+\varphi(x_1,x_2)+\varepsilon(x_3).
\]
A further direct calculation gives
\[
[T,H_0-b_0D_2]=\frac12\ell_1'(x_3)
\]
and
\[
[T,H_1+b_0D_1]
=\frac12\bigl(\ell_2'(x_3)-2h_0\varepsilon'(x_3)\bigr).
\]
Both sides are function elements depending only on \(x_3\), and hence
are constants. Thus
\[
\ell(x_3)
=2h_0\varepsilon(x_3)
\begin{pmatrix}0\\1\end{pmatrix}
+a+x_3a'
\]
for some constant vectors \(a,a'\in\mathbb R^2\). We have proved
\begin{equation}
\eta
=\eta_0(x_1,x_2)
+2\varepsilon(x_3)F(x_1,x_2)
+e(x_3)
+(a+x_3a')\cdot y.
\label{eq:sectorII-eta}
\end{equation}

\medskip
\noindent\emph{Step 2: a hidden gauge and polynomiality.}

Choose a smooth univariate function \(\Lambda\) satisfying
\[
\Lambda'=\varepsilon
\]
and set
\[
\widetilde E=e^\Lambda Ee^{-\Lambda}.
\]
This conjugation fixes multiplication operators and sends
\[
D_3=\partial_3-F
\quad\text{to}\quad
\partial_3-(F+\varepsilon).
\]
It also sends \(T\) to
\[
\widetilde T=\partial_3+g(x_1,x_2)
\]
for some smooth function \(g\).

Set
\[
r(x_3):=e(x_3)-\varepsilon(x_3)^2.
\]
Expanding the conjugated \(L_0\) and using
\eqref{eq:sectorII-eta}, the terms
\(2\varepsilon F+\varepsilon^2\) cancel. Hence
\begin{align}
\widetilde L_0={}&
\frac12\bigl(\partial_1^2+\partial_2^2+\partial_3^2\bigr)
-b_0x_1\partial_2-(F+\varepsilon)\partial_3\notag\\
&+V_0(x_1,x_2)+x_3V_1(x_1,x_2)
-\frac12\bigl(\varepsilon'(x_3)+r(x_3)\bigr),
\label{eq:sectorII-L}
\end{align}
where the precise forms of the visible functions \(V_0,V_1\) will not
be needed.

The coefficient of \(\partial_3\) in
\[
[\widetilde L_0,\widetilde T]
\]
is \(\varepsilon'(x_3)\). If this commutator has order one, the theorem
on polynomial highest-order coefficients, transported through the gauge
conjugation, shows that \(\varepsilon'\) is a polynomial. If it has order
zero, then \(\varepsilon'=0\). Thus
\[
\varepsilon\in\mathbb R[x_3].
\]

Let \(d=\deg\varepsilon\), with \(d=0\) when \(\varepsilon=0\), and put
\[
m=\max\{2,d\},
\qquad
c=\varepsilon^{(m)}.
\]
Then \(c\) is a real constant. Since
\[
\widetilde T=\partial_3+g(x_1,x_2),
\qquad
[\partial_3,g]=0,
\]
and every commutator with the multiplication operator \(g\) lowers
differential order, direct expansion gives
\[
\operatorname{ad}_{\widetilde T}^{\,m}(\widetilde L_0)
+c\widetilde T
\]
as a multiplication operator. Its only \(x_3\)-dependent part is
\[
-\frac12r^{(m)}(x_3).
\]
It belongs to
\[
\widetilde E\cap C^\infty(\mathbb R^3)
=\operatorname{span}_{\mathbb R}\{1,x_1,x_2\},
\]
so \(r^{(m)}\) is constant. Therefore
\[
r\in\mathbb R[x_3],
\qquad
\deg r\le\max\{2,\deg\varepsilon\}.
\]

\medskip
\noindent\emph{Step 3: a weighted Lie ladder.}

Every coefficient in \eqref{eq:sectorII-L} is now polynomial in \(x_3\).
For a normally ordered monomial define
\[
\nu\!\left(
a(x_1,x_2)x_3^p
\partial_1^{q_1}\partial_2^{q_2}\partial_3^{q_3}
\right)
:=p+q_1+q_2+2q_3.
\]
The visible coefficient \(a(x_1,x_2)\) is not counted. The weight of an
operator is the largest weight of its monomials.

A direct application of the Leibniz rule gives the only estimates needed:

\begin{itemize}
\item the constant second-order part and \(-F\partial_3\) raise the
weight by at most \(1\);
\item \(-c_sx_3^s\partial_3\) raises it by at most \(s-1\);
\item multiplication by a polynomial of degree \(D\) raises it by at
most \(D-3\);
\item all remaining terms in \eqref{eq:sectorII-L} raise it by at most
\(0\).
\end{itemize}

Consider
\[
S_n
:=\bigl(\operatorname{ad}_{\widetilde L_0}\bigr)^n(\partial_1) \in \widetilde E\]
(i) If $\deg \eps \ge 3$:

Write
\[
\varepsilon(x_3)
=ax_3^d+\text{lower powers},
\qquad a\ne0.
\]
Since \(\deg r\le d\), the unique term that can raise the weight by
\(d-1\) is
\[
-ax_3^d\partial_3.
\]
Moreover,
\[
S_1
=kx_1\partial_3+\text{terms of weight at most \(1\)}
\]
and
\[
[-ax_3^d\partial_3,\,
  x_1x_3^q\partial_3]
=
a(d-q)x_1x_3^{q+d-1}\partial_3.
\]
It follows inductively that the highest-weight term of \(S_{n+1}\) is
\[
ka^n
\prod_{j=0}^{n-1}\bigl(d-j(d-1)\bigr)
x_1x_3^{n(d-1)}\partial_3.
\]
with $d - j(d-1) \neq 0$ for $d \ge 3$ as an integer. This derives contradiction.

(ii): \(\deg\varepsilon\le2\), including
\(\varepsilon=0\).

Let \(a\) be the coefficient of \(x_3^2\) in \(\varepsilon\), allowing
\(a=0\). Every bracket with \(\widetilde L_0\) now raises the weight by
at most \(1\).

Let \(\xi,\zeta\) be the graded symbols of
\(\partial_1,\partial_3\). Starting from \(\xi\), the part that raises
the weight by exactly \(1\) is the derivation
\[
\mathcal D
=
\xi\partial_{x_1}
+(\zeta-ax_3^2)\partial_{x_3}
+kx_1\zeta\,\partial_\xi
+2ax_3\zeta\,\partial_\zeta.
\]
The terms containing \(h\) and \(h_0\) act only on the
\((x_2,\partial_2)\)-variables and therefore never enter
\(\mathcal D^n\xi\).

We claim that
\[
\mathcal D^n\xi\ne0
\qquad(n\ge0).
\]
Otherwise, suppose that \(\mathcal D^N\xi=0\). Along every local integral
curve of \(\mathcal D\), both \(\xi(\tau)\) and \(x_1(\tau)\) would then
be polynomials. The relevant equations are
\[
\dot x_1=\xi,\qquad
\dot\xi=kx_1\zeta,\qquad
\dot x_3=\zeta-ax_3^2,\qquad
\dot\zeta=2ax_3\zeta.
\]
Choose
\[
x_3(0)=0,\qquad \zeta(0)=1.
\]
For this fixed hidden trajectory, the two initial conditions
\[
(x_1(0),\xi(0))=(1,0),
\qquad
(x_1(0),\xi(0))=(0,1)
\]
give two polynomial solutions \(p,q\) of
\[
u''=k\zeta(\tau)u.
\]
Their Wronskian is the constant \(1\).

Two polynomials with a nonzero constant Wronskian span a nonzero constant
polynomial: after making their degrees unequal by a constant change of
basis, the degree of their Wronskian is the sum of their degrees minus
one. Hence those degrees must be \(0\) and \(1\). The equation therefore
has a nonzero constant solution \(u\), which would imply
\[
0=k\zeta(\tau)u.
\]
This contradicts \(k\ne0\) and \(\zeta(0)=1\). Thus
\[
\mathcal D^n\xi\ne0
\qquad(n\ge0).
\]
Consequently, the highest-weight term of \(S_n\) is
\(\mathcal D^n\xi\), and
\[
\nu(S_n)=n+1\longrightarrow\infty.
\]


Therefore \(M=0\). Returning to the original visible coordinates gives
\[
k_1=k_2=h_2=0.
\]
The proof divides only by \(k\ne0\), so it includes the rank-one case
\(h=0\) and allows arbitrary \(h_0\).
\end{proof}


\section{Conclusion}

Every unbounded sequence in the proof consists of actual Lie words.  The
proof for the visible entry uses only the function-element space and the
affine Wong structure.  In Sector I, the spectral projections are introduced
only after Cayley--Hamilton and the global smooth Euler equation have produced
a finite spectrum, and they are then realized by the Chinese remainder
theorem.  Sector II is treated inside
$C^\infty(\R^2)[x_3]\langle\pd_1,\pd_2,\pd_3\rangle$.  Thus we never treat
the Lie algebra as a module over functions, nor do we use analytic
continuation, maximal rank, or an a priori polynomial $\eta$.

Jiao's Theorem~5.4 also derives, from the structure of $\eta$, polynomiality
of the pre-gauge term $r$.  The smooth Euler lemma in the present paper
additionally controls the post-gauge zeroth-order slot $U$, so the final
argument remains strictly within the smooth assumptions of the original
problem.

\begin{thebibliography}{99}

\bibitem{Ocone1981}
D. L. Ocone,
\newblock Finite dimensional estimation algebras in nonlinear filtering,
\newblock in M. Hazewinkel and J. C. Willems (eds.),
\emph{Stochastic Systems: The Mathematics of Filtering and Identification
and Applications}, D. Reidel, Dordrecht, 1981, pp. 629--636,
\href{https://doi.org/10.1007/978-94-009-8546-9_35}
{doi:10.1007/978-94-009-8546-9\_35}.

\bibitem{ShiYau2017}
J. Shi and S. S.-T. Yau,
\newblock Finite dimensional estimation algebras with state dimension 3
and rank 2, I: linear structure of Wong matrix,
\newblock \emph{SIAM Journal on Control and Optimization}
55 (2017), no. 6, 4227--4246,
\href{https://doi.org/10.1137/16M1065471}{doi:10.1137/16M1065471}.

\bibitem{ShiYau2020}
J. Shi and S. S.-T. Yau,
\newblock Finite dimensional estimation algebras with state dimension 3
and rank 2, Mitter conjecture,
\newblock \emph{International Journal of Control}
93 (2020), no. 9, 2177--2186,
\href{https://doi.org/10.1080/00207179.2018.1550268}
{doi:10.1080/00207179.2018.1550268}.

\bibitem{Jiao2022}
Xiaopei Jiao,
\newblock \emph{The Classification of Finite Dimensional Estimation Algebra
and New Classes of Finite Dimensional Filters},
\newblock Ph.D. thesis in Mathematics, Tsinghua University, 2022
(in Chinese).

\end{thebibliography}

\end{document}
 
